\section{Encoding and continuum geometry}\label{sec:framework}

In the FK proof \(q\in(0,4)\) is fixed. Let \(d_\gamma\) be the
almost sure Hausdorff dimension of the \(\gamma\)-LQG metric,
a deterministic constant \cite[Corollary~1.7]{GwynnePfefferKPZ}. Put
\begin{equation}\label{bg:parameters}
 b_{\mathrm F}=\frac{\sqrt q}{2+\sqrt q},\qquad
 \rho=\frac{b_{\mathrm F}}{1-b_{\mathrm F}},\qquad
 \theta=\frac{\pi\gamma^2}{4},\qquad
 p=\frac1{d_\gamma},\quad \xi=\gamma p,\quad
 Q=\frac2\gamma+\frac\gamma2 .
\end{equation}
Thus \(\rho=-\cos\theta\), \(\pi/2<\theta<\pi\), and
\begin{equation}\label{bg:strict-margins}
 Q-2=\frac{(\gamma-2)^2}{2\gamma}>0,\qquad
 \frac{\pi}{2\theta}-\frac12=\frac2{\gamma^2}-\frac12>0.
\end{equation}
All constants may depend on \(q\). A constant depending on \(q\) is fixed before any scale tends to infinity.

We first identify the word law and its continuum surface, then record the
local field and metric properties used in the proof. These inputs identify
the limiting geometry of the encoding; they do not assert convergence of
discrete graph distances.

\subsection{The discrete and continuum encodings}

An inventory letter is a burger of type \(1\) or \(2\), a rigid order of one of these types, or a flexible order. Their respective probabilities are
\[
 \frac14,\quad\frac14,\quad
 \frac{1-b_{\mathrm F}}4,\quad\frac{1-b_{\mathrm F}}4,
 \quad\frac{b_{\mathrm F}}2.
\]
An order takes the freshest available burger of an allowed type. Reduction commutes orders past burgers of the other type and cancels a burger with its matching order. We use both a bilateral iid word and a word of length \(2n\) conditioned to reduce to the empty word. After the flexible orders are resolved, the two net burger counts give the contour \(Z=(L,R)\).

We use the inventory bijection in the convention of \cite{SheffieldInventory}. It produces the primal all-edge map, its tree/cotree sewing, its FK decoration, and its oriented root. The primal tour vertex at integer time \(j\) is denoted by \(v_n(j)\). Consecutive tour vertices are equal or joined by a primal tree edge. A match in the other contour supplies a remaining primal edge through its quadrangle. Consequently every graph edge used below is an edge of the primal map, whether or not it belongs to the FK subgraph.

\begin{lemma}[The exact finite law]\label{bg:finite-law}
The inventory word of length \(2n\) conditioned to reduce to the empty word induces the joint law proportional to \(q^{\ell(M,A)/2}\) in the statement of \Cref{thm:main}. Conditional on the unrooted decorated map, its root corner is uniform.
\end{lemma}
\begin{proof}
An empty word has \(n\) burger letters and \(n\) order letters. If \(f\) of its orders are flexible, its unconditioned probability is
\[
 4^{-n}\left(\frac{1-b_{\mathrm F}}4\right)^{n-f}
       \left(\frac{b_{\mathrm F}}2\right)^f
 =c_{n,q}\left(\frac{2b_{\mathrm F}}{1-b_{\mathrm F}}\right)^f.
\]
The bijection has \(\ell=f+1\). Since \(2b_{\mathrm F}/(1-b_{\mathrm F})=\sqrt q\), the displayed weight is a factor independent of \((M,A)\) times \(q^{\ell/2}\). Rooted isomorphism classes correspond under the same bijection. Re-rooting does not change the weight. Equivalently, on an unrooted class one may choose a uniform dart; the automorphism group acts freely on darts, so this gives the asserted rooted law. A corner and its outgoing dart are equivalent root conventions.
\end{proof}

The public continuum inputs will be used in the following precise forms.

\begin{theorem}[Contour and surface inputs]\label{bg:continuum-inputs}
After fixed deterministic changes of time and height units:
\begin{enumerate}[label=(\roman*)]
\item The bilateral resolved contour, viewed at time scale \(m\) and height scale \(\sqrt m\), converges locally uniformly in law to a two-dimensional Brownian motion with equal coordinate variances and correlation \(\rho\). The empty-word contour at total time \(2n\), with total rescaled time one, converges to the corresponding quadrant excursion from zero to zero.
\item The bilateral Brownian contour determines a curve-decorated \(\gamma\)-quantum cone. Its curve \(\eta:\R\to\C\) is continuous, transient, space filling, and parameterized by quantum area. The excursion determines the ordinary unit-area DMS quantum sphere with its area-parameterized space-filling curve. The latter sphere is the quantum-sphere measure disintegrated at area one.
\item The spatial equivalence relation of the sewing is generated by finite chains of scalar horizontal chords. For a scalar coordinate \(X\), a horizontal chord joins \(s<t\) if
\[
 X(s)=X(t)=\inf_{[s,t]}X.
\]
At the deterministic or independently sampled cuts used below, the two boundary strands are contour-tree lineages parameterized by quantum boundary length. In the cone coupling, the intrinsic decorated surface of each fixed bounded time interval is determined by its contour increments. This assertion may be used simultaneously for a countable collection of deterministic cuts or cuts sampled independently of the surface, and for later selection among that collection.
\item On a buffered spatial patch, the imaginary field determines its directed local traversals and their frontier flow lines. Together with the real field and local Poisson area marks, these determine the clocked local pieces, their incoming and outgoing sides, and the quantum lengths of the local frontier arcs. This statement does not determine the global order of separate traversals.
\end{enumerate}
\end{theorem}
\begin{proof}[References and conventions]
Part (i) is the bilateral inventory invariance principle of \cite{SheffieldInventory} and the finite-volume result \cite[Theorem~1.8]{GwynneSunFinite}. Parts (ii)--(iii) are the mating-of-trees results of \cite{DMS} and \cite[Theorem~1.1]{MillerSheffieldSphere}; cone interval-surface determination in the beaded regime is stated explicitly in \cite[Lemma~3.13]{GwynneHoldenSunBipolar}. Independent random cuts follow by conditioning on their values, and countable intersections allow the stated later selections. This is intrinsic determination, modulo quantum-surface equivalence. The finite-sphere proof below instead transfers spatial observations using actual local fields; it does not require a sphere version of the isolated-interval reconstruction statement.

The finite-chain convention in (iii) is the strict-chord relation of \cite[Lemma~8.9]{DMS}, with its generated classes identified in \cite[Lemmas~8.14--8.15]{DMS}. A weak scalar chord splits into at most two strict chords: otherwise two distinct interior local minima would have the same height, excluded by \Cref{arm:generic-chords}. The same frontier-tree relation applies to the sphere's whole-plane space-filling curve. Quantum-area reparametrization and its quantum boundary-length coordinates preserve the tree identifications; the initial and final times are identified at the exploration root.

For (iv), \cite[Lemma~2.4]{GwynneMillerSheffieldHarmonic} reconstructs directed parameterized local traversals with a buffer; its quantum-area clock is described in the proof of \cite[Lemma~4.3]{GwynneMillerSheffieldHarmonic}. The frontier arcs are the labelled local flow-line boundaries in that construction. Their quantum lengths use the boundary Gaussian multiplicative chaos construction and its conformal coordinate rule \cite[Sections~1.2 and~5.1]{SheffieldQuantumZipper}. More explicitly, choose the specified side of a compact frontier subarc and a sided domain inside the buffer, then straighten that domain conformally and construct boundary length from the restriction of the actual normalized real field. The boundary coordinate rule identifies this measure with the same side's global quantum boundary length and makes it independent of the chosen smaller domain. Exhaust the arc by compact subarcs to include its ends. The sided domains and prime-end labels come from the local imaginary field, so this construction uses only the indicated field restrictions. Contained Poisson-cell locality is also given in \cite[Lemma~3.7]{ContrerasHipGwynneCurvature}. We always shrink the testing region before using these locality statements.

The imaginary whole-plane field is taken modulo its angular period, with its stationary phase, rather than circle-average-pinned with a prescribed phase. Local decompositions consist of Dirichlet GFF parts and harmonic parts modulo that period. The real and imaginary fields are independent before the curve is clocked by real-field area.
\end{proof}

\subsection{Ordinary patches of the sphere}

The finite-volume argument will transfer local observations from a quantum
cone to the ordinary sphere. Such a comparison must retain the field's
additive constant: that constant fixes both area and intrinsic distance.
The next lemma proves the required one-way domination with the actual
unit-area normalization in place.

\begin{lemma}[Actual-field absolute continuity on a sphere patch]\label{bg:sphere-field-ac}
Let \((\mathbb C,h,\infty,0,1)\) be the ordinary unit-area quantum
sphere with three conditionally independent quantum-area marks in the
indicated embedding. If \(U\) is bounded and
\(\overline U\subset\mathbb C\setminus\{0,1\}\), and \(R\) is large
enough that a neighborhood of \(\overline U\) lies in \(B_R(0)\), then
\[
 \operatorname{Law}(h|_U)
 \ll \operatorname{Law}(G_R|_U),
\]
where \(G_R\) is the whole-plane GFF pinned to have zero average on
\(\partial B_R(0)\). This is a statement about actual distributions,
with their additive constants retained. The right-hand law is mutually
absolutely continuous with the restriction of a deterministically
dilated circle-average embedding of a \(\gamma\)-quantum cone whose
normalization circle is \(\partial B_R(0)\). Both comparisons remain valid
jointly with the complete independent stationary-phase imaginary field
and its associated unparameterized whole-plane space-filling curve.
\end{lemma}

\begin{proof}
We use \cite[Definition~2.2]{BorgaGwynneSunPermutons}, whose
three-marked construction agrees with the ordinary DMS sphere by
\cite[Theorem~1.1 and Remark~1.2]{AruHuangSunSphere}.
Write \(G_1\) for the whole-plane GFF pinned on the unit circle,
and let \(d\) be the deterministic Liouville drift in that definition.
The sphere field is obtained by a positive finite integrable weighting
of
\[
 N(G_1+d),\qquad
 N(H)=H-\gamma^{-1}\log\mu_H(\mathbb C).
\]
The drift has logarithmic singularities only at \(0,1,\infty\);
on compact subsets avoiding \(0,1\) it has finite Dirichlet energy.
In particular, the possible crease of \(\log(|z|\vee1)\) at the unit
circle causes no difficulty for Cameron--Martin comparison.

Couple the pinned GFFs by
\(G_1=G_R-(G_R)_1(0)\). Since \(N(H+c)=N(H)\), including for a
field-dependent scalar \(c\), the normalized field is also
\(N(G_R+d)\). The original weighting is still a positive finite
integrable density on this Gaussian probability space, although its
formula after changing the pin need not be the original one. It
therefore suffices first to prove the domination for the unweighted
field \(N(G_R+d)\).

Choose \(g\in C_c^\infty(B_R(0))\) with \(0\le g\le1\), equal to one
on a neighborhood of \(\overline U\). Gaussian orthogonal
decomposition in the Cameron--Martin direction \(g\) gives
\(G_R=Xg+G^\perp\), where \(X\) is a nondegenerate real Gaussian
independent of \(G^\perp\). Conditional on \(G^\perp\), set
\(\nu=\mu_{G^\perp+d}\) and
\[
 A(x)=\int_{\mathbb C}e^{\gamma xg(z)}\,\nu(dz),
 \qquad T(x)=x-\gamma^{-1}\log A(x).
\]
Almost surely \(\nu\) is finite, nonzero, and assigns positive mass to
an open set outside \(\operatorname{supp}g\). These assertions follow
from the corresponding properties of \(\mu_{G_R+d}\), since removing
the bounded function \(Xg\) changes that measure by a bounded positive
factor. The same comparison applies to every finite value of \(x\).
Differentiation under the integral on compact \(x\)-intervals yields
\[
 T'(x)=
 \frac{\displaystyle\int_{\mathbb C}(1-g(z))e^{\gamma xg(z)}\,\nu(dz)}
      {\displaystyle\int_{\mathbb C}e^{\gamma xg(z)}\,\nu(dz)}>0 .
\]
Thus \(T\) is a strictly increasing \(C^1\) diffeomorphism onto its
image. Conditional on \(G^\perp\), the law of \(T(X)\) has a
Lebesgue density and hence is absolutely continuous with respect to
the nondegenerate Gaussian law of \(X\). On \(U\),
\[
 N(G_R+d)|_U=G^\perp|_U+d|_U+T(X).
\]
Integrating the preceding conditional domination proves
\[
 \operatorname{Law}(N(G_R+d)|_U)
 \ll \operatorname{Law}((G_R+d)|_U).
\]
A finite-energy cutoff equal to \(d\) near \(\overline U\), supported
away from \(0,1\) and the pinning circle, is a Cameron--Martin shift
of \(G_R\). Its restriction removes \(d|_U\) by mutual absolute
continuity. Restoring the positive weighting proves the first assertion.

For the comparison with the cone, \cite[Definition~4.10 and the
following discussion]{DMS} identifies its field inside the unit disk
with \(G_1-\gamma\log|z|\). Under the deterministic coordinate map
\(z\mapsto z/R\), its representation on \(B_R(0)\) is
\[
 G_R-\gamma\log|z|+(\gamma-Q)\log R.
\]
On \(\overline U\), the added deterministic function has a
finite-energy cutoff away from \(0\) and \(\partial B_R(0)\).
Cameron--Martin gives mutual absolute continuity with \(G_R|_U\).
In particular the direction needed for transferring a cone
probability-one event is
\[
 \operatorname{Law}(h|_U)
 \ll \operatorname{Law}(G_R|_U)
 \ll \operatorname{Law}(h^{\mathrm{cone},R}|_U).
\]
No comparison across the cone's pinning circle is required.

Finally, the complete imaginary field is independent of the real
field and has the same stationary-phase law in both constructions.
Taking the product with this common law, and then taking measurable
images to obtain its curve, preserves each asserted domination.
\end{proof}

\begin{remark}[Auxiliary marks and affine charts]\label{bg:sphere-affine-marks}
In a singly marked sphere with exploration root \(p_*=\infty\),
sample two auxiliary points \(Q_1,Q_2\) independently from quantum
area, conditionally on the surface and independently of the imaginary
field. Sending them to \(0,1\) by the unique complex affine map gives
the three-marked embedding above, by
\cite[Lemma~2.3]{BorgaGwynneSunPermutons} and
\cite[Proposition~A.13]{DMS}. Explicitly, if
\(\psi(z)=Q_1+(Q_2-Q_1)z\) maps the new coordinate to the old one,
then \(h_{\mathrm{new}}=h\circ\psi+Q\log|\psi'|\).
The affine map is independent of the imaginary field; conditional on
the real surface and marks, affine invariance of the whole-plane
space-filling SLE and the stationary imaginary-field convention
preserve its complete law. Thus the transformed imaginary field
remains independent of the transformed real field, before quantum-area
parametrization.

A second independent auxiliary pair \(Q_3,Q_4\) gives another such
chart with the same exploration root. The two pairs are disjoint
almost surely. Every point other than \(p_*\) is ordinary in at least
one chart, so countably many compactly buffered patches in the two
charts cover the punctured sphere. This is a local covering assertion;
it does not assert that every large curve interval lies in one
ordinary patch.
\end{remark}

\subsection{The reference intrinsic metric}

The metric below is the already constructed LQG metric. Later we compare
an extracted discrete passage rule with it. At present that passage rule
has neither an established metric property nor a Weyl-scaling law.

\begin{theorem}[Metric inputs]\label{bg:metric-inputs}
For the fixed subcritical parameter, the LQG metric \(D_h\) is a continuous length metric inducing the Euclidean topology. In whole-plane GFF coordinates it is proper and geodesic, and geodesics between each fixed pair of distinct points are almost surely unique. Internal metrics are local in the field. Conformal covariance and Weyl scaling hold: adding a continuous function \(f\) multiplies path length locally by \(e^{\xi f}\).

For a compact set \(K\) in an ordinary field chart and
\(0<\chi<\xi(Q-2)\), there are almost surely finite random constants controlling connections between nearby points of \(K\) inside Euclidean balls of a fixed multiple of their separation by \(C|z-w|^\chi\). The corresponding internal diameters of small dyadic squares, with a bounded enlargement still in the chart, have the same type of bound. Fixed boxes of bounded aspect ratio admit access to their flat faces by completion of the internal metric.

In whole-plane GFF coordinates, strong confluence holds simultaneously for all geodesics: an interior segment of any geodesic persists in all geodesics whose endpoints sufficiently closely approximate slightly longer interior subsegments. The same local assertion holds for short geodesic subarcs in buffered ordinary surface charts. In particular, every nonempty open geodesic subarc contains a segment of a geodesic between points of a deterministic countable dense set.
\end{theorem}
\begin{proof}[References and use]
Metric construction, locality, and Weyl scaling are supplied by \cite{GwynneMillerMetric}; properness is \cite[Lemma~3.8]{DFGPSWeakMetric}, and general conformal coordinate changes are \cite[Theorem~1.3]{GwynneMillerConformal}. Fixed-pair geodesic uniqueness is \cite[Theorem~1.2]{MillerQianGeodesics}. The internal connection and dyadic-square estimates are \cite[Lemma~3.20, Equations~(3.66)--(3.67)]{DFGPSWeakMetric}. Strong confluence and its interior-endpoint form are \cite[Theorem~2 and Proposition~6]{BhatiaKavvadiasConfluence}. Local statements transfer to fixed compact subsets of ordinary cone or sphere charts by local absolute continuity. For confluence, first take a shorter geodesic subarc in a buffered chart. Its sufficiently short neighboring connections cannot reach the chart boundary, whose internal distance from a compact inner neighborhood is positive. They are consequently determined by the same internal metric, so the whole-plane confluence event applies locally. This does not require absolute continuity of an entire surface's global geodesics. Singular marked points will be treated separately.

For clarity, the open-square internal bound in \cite[Proposition~3.10]{DFGPSWeakMetric} does not by itself include boundary points. To access a specified flat-face or corner point of a rectangle of size \(\delta\), choose inward points at Euclidean distances comparable to \(2^{-j}\delta\). Consecutive inward points can be joined inside the rectangle by a bounded chain of balls or dyadic squares of that scale. The internal connection bound gives total cost at most
\[
 C\sum_{j\ge0}(2^{-j}\delta)^\chi
 =\frac{C\delta^\chi}{1-2^{-\chi}}.
\]
The chain is Cauchy for the internal metric. Its Euclidean limit is the prescribed boundary point, and the same inward-chain construction between two approximations shows uniqueness in the internal completion. A bounded number of these connections also crosses any fixed-shape polygonal collar while staying inside that collar. No bound uniform over degenerating aspect ratios is claimed.

These properties are asserted only for \(D_h\). No analogous Weyl rule, metric property, or endpoint modulus is assumed for a subsequential discrete passage rule.
\end{proof}
